% intro.tex

%%%%%%%%%%%%%%%%%%%%
\begin{frame}{}
  \begin{center}
    \blue{\bf (数理) 逻辑 (Mathematical Logic)} \\[3pt]
    是研究 \teal{\bf 证明/推理 (Proof/Inference)} 的一门学科。

    \vspace{0.30cm}
    \fig{width = 0.60\textwidth}{figs/euler-logic}

    \pause
    \vspace{0.30cm}
    \red{\bf 一个基本问题: 什么样的推理(规则)是正确/合法 (Valid) 的?}

    \pause
    \vspace{0.50cm}
    将\red{\bf ``证明/推理''}作为数学对象, 研究其结构、性质、\cyan{局限性}等问题。
  \end{center}
\end{frame}
%%%%%%%%%%%%%%%%%%%%

%%%%%%%%%%%%%%%%%%%%
\begin{frame}{}
  \begin{center}
    \fig{width = 0.60\textwidth}{figs/math-tree}

    \vspace{0.50cm}
    逻辑要研究的是\red{\bf 不依赖于具体领域知识}的\teal{\bf 推理}
  \end{center}
\end{frame}
%%%%%%%%%%%%%%%%%%%%

%%%%%%%%%%%%%%%%%%%%
\begin{frame}{}
  \[
    \ndnoname{P \text{ 是真的} \qquad {Q \text{ 是真的}}}{P \text{ 与}\; Q \text{ 都是真的}}
    \qquad \ndnoname{P \qquad Q}{P \;\red{\land}\; Q}
  \]

  \pause
  \[
    \ndnoname{\text{如果 $P$ 是真的, 则 $Q$ 是真的} \qquad {\text{$P$ 是真的}}}{\text{$Q$ 是真的}}
    \qquad \ndnoname{P \;\red{\to}\; Q \qquad P}{Q}
  \]

  \pause
  \[
    \gray{\ndnoname{\text{如果 $P$ 是真的, 则 $Q$ 是真的} \qquad {\text{$Q$ 是真的}}}{\text{$P$ 是真的}}
    \qquad \ndnoname{P \;\red{\to}\; Q \qquad Q}{P}}
  \]

  \pause
  \[
    \ndnoname{\text{如果 $P$ 是真的, 则 $Q$ 是真的} \qquad {\text{$Q$ 是假的}}}{\text{$P$ 是假的}}
    \qquad \ndnoname{P \to Q \qquad \red{\lnot}\; Q}{\red{\lnot}\; P}
  \]

  \pause
  \begin{columns}
    \column{0.60\textwidth}
      \begin{center}
        \blue{这些都是人类最基本的思维规律}
      \end{center}
    \column{0.40\textwidth}
      \fig{width = 0.40\textwidth}{figs/Theophrastus}
  \end{columns}
\end{frame}
%%%%%%%%%%%%%%%%%%%%

%%%%%%%%%%%%%%%%%%%%
\begin{frame}{}
  \begin{center}
    \fig{width = 0.40\textwidth}{figs/logic-logo}

    \vspace{0.80cm}
    逻辑就是\red{\bf 用数学的方法}研究\blue{\bf 人类基本思维规律}的一门学科
  \end{center}
\end{frame}
%%%%%%%%%%%%%%%%%%%%

%%%%%%%%%%%%%%%%%%%%
\begin{frame}{}
  \begin{center}
    \fig{width = 0.60\textwidth}{figs/prop-logic-vs-predicate-logic}

    \vspace{0.80cm}
    先来学习最基础的\blue{\bf 命题逻辑 (Propositional Logic)}
  \end{center}
\end{frame}
%%%%%%%%%%%%%%%%%%%%

%%%%%%%%%%%%%%%%%%%%
\begin{frame}{}
  \begin{center}
    命题逻辑是一个\red{\bf 公理系统}, \\[3pt]
    包含 \teal{\bf 语法 (Syntax)} 与 \teal{\bf 语义 (Semantics)} 两部分。

    \vspace{0.30cm}
    \fig{width = 0.60\textwidth}{figs/syntax-semantics}
    \vspace{0.30cm}

    \pause
    语法与语义是\red{\bf ``对立统一''}的

    \pause
    \vspace{0.50cm}
    {\large \blue{\bf 螺旋上升:} 先讲语法, 再讲语义, 再回到语法, 最后将二者统一起来}
  \end{center}
\end{frame}
%%%%%%%%%%%%%%%%%%%%

% %%%%%%%%%%%%%%%%%%%%
% \begin{frame}{}
%   \[
%     \red{\set{\alpha, \beta, \gamma, \dots}} \;\cyan{\models}\; \blue{\tau}
%   \]

%   \vspace{0.50cm}
%   \begin{center}
%     \cyan{\bf 正确性:} 如果\red{\bf 前提} $\alpha$, $\beta$, $\gamma$, $\dots$ 都为真,
%     那么\blue{\bf 结论} $\tau$ 也为真。
%   \end{center}
% \end{frame}
% %%%%%%%%%%%%%%%%%%%%

% %%%%%%%%%%%%%%%%%%%%
% \begin{frame}{}
%   \begin{center}
%     证明就是一系列根据\cyan{\bf 推理规则}将公式不断变形的步骤。

%     \[
%       \red{\set{\alpha, \beta, \gamma, \dots}} \;\cyan{\vdash}\; \blue{\tau}
%     \]

%     \[
%       \ndnoname{\alpha \qquad \beta \qquad \gamma}{\tau}
%     \]

%     \pause
%     \vspace{0.50cm}
%     每一步推理都必须是正确的

%     \pause
%     \vspace{0.60cm}
%     \red{\bf 什么样的推理是正确/合法 (Valid) 的?}
%   \end{center}
% \end{frame}
% %%%%%%%%%%%%%%%%%%%%